\documentclass[11pt]{article}

\usepackage[utf8]{inputenc}
\usepackage[T1]{fontenc}
\usepackage{kotex}
\usepackage{lmodern}
\usepackage{amsmath}
\usepackage{amssymb}
\usepackage{amsthm}
\usepackage{mathtools}
\usepackage{bm}
\usepackage{geometry}
\usepackage{enumitem}
\usepackage{array}
\usepackage{booktabs}
\usepackage{hyperref}

\geometry{margin=1in}
\hypersetup{
  colorlinks=true,
  linkcolor=blue,
  citecolor=blue,
  urlcolor=blue
}

\newtheorem{definition}{정의}
\newtheorem{axiom}{공리}
\newtheorem{proposition}{명제}
\newtheorem{assumption}{가정}
\theoremstyle{remark}
\newtheorem{remark}{비고}

\renewcommand{\abstractname}{초록}
\renewcommand{\contentsname}{목차}
\renewcommand{\refname}{참고문헌}

\title{Jaemin Cosmology: A Higher-Dimensional Information Projection Framework for Consciousness, Intuition, and Internal Time}
\author{Jaemin}
\date{\today}

\begin{document}

\maketitle

\begin{abstract}
본 논문은 ``Jaemin Cosmology'', 즉 ``재민 우주론''이라는 이름의 사변적 수학
틀을 형식화한다. 초기 동기는 내부 시간 가속과 주관적 시간 변화에 관한
의문이지만, 본 논문이 정식화하는 중심 명제는 뇌가 단순한 삼차원 계산 장치가
아니라 고차원 정보공간과 삼차원 의식 경험을 연결하는 인터페이스일 수
있다는 가설이다. 이 틀에서 인간이 경험하는 삼차원 공간은 현실 전체가 아니라
인지적 투영이며, 의식은 고차원 정보 현실의 저차원 인터페이스, 잠재의식은
더 높은 차원의 표현, 직관은 그러한 표현이 의식 수준으로 압축된 신호로
정의된다. 또한 뇌의 연결 복잡도와 고차원 정보 접근 효율을 구별하여 지능,
통찰, 창의성의 차이를 해석하는 수학적 가능성을 제시한다. 내부 우주 시간과
고차원 기준 시간을 구별함으로써, 고차원 관찰자에게는 내부 시간이 가속되는
것처럼 보이지만 내부 관찰자는 자신의 물리 법칙을 정상적으로 측정할 수
있음을 보인다. 그러나 이 글은 검증된 물리 이론이나 신경과학 이론을
제시하지 않는다. 현재 단계의 재민 우주론은 경험적으로 확인된 과학 이론이
아니라, 수학적으로 정의 가능한 해석적 가설이자 이론적 프레임워크이다.
\end{abstract}

\tableofcontents

\section{서론 (Introduction)}

인간의 의식 경험은 삼차원 공간과 시간의 흐름으로 구성된 세계를 기본
형식으로 받아들인다. 그러나 어떤 지각 체계가 경험하는 차원 수가 곧 현실의
전체 차원 수라는 결론은 논리적으로 필연적이지 않다. 이 논문은 인간
의식이 고차원 정보 구조의 제한된 저차원 인터페이스일 수 있다는 가설을
수학적으로 정리한다.

이 프레임워크의 출발점은 시간 가속에 관한 사고실험이다. 만약 내부 우주의
모든 물리 과정이 외부 기준 시간에 대해 동일한 비율로 빨라진다면, 내부
관찰자는 그 변화를 직접 감지하지 못한다. 그러나 이 사고실험은 본 논문의
중심 결론이 아니라 동기이다. 논의를 확장하면 더 중요한 질문은 인간이
어떻게 직관, 통찰, 창의성, 지능 차이를 경험하는가이다. 재민 우주론의 핵심
명제는 다음과 같다.

\begin{quote}
뇌는 현실을 완전히 내부에서 계산하는 장치가 아니라, 고차원 정보공간과
삼차원 의식 인터페이스 사이의 변환기일 수 있다.
\end{quote}

이 논문의 목적은 물리학의 표준 모형이나 일반상대성이론을 대체하는 것이
아니다. 또한 경험적 확증을 주장하지 않는다. 본문의 모든 정의와 명제는
다음 세 수준을 구별한다.

\begin{description}[leftmargin=2.8cm, style=nextline]
  \item[Level A: 수학적 가능성]
  이 프레임워크는 명시적 사상, 상태공간, 시간 매개변수, 정보량 조건으로
  정의될 수 있다.

  \item[Level B: 해석적 이론]
  의식, 잠재의식, 직관, 지능 차이, 주관적 시간 경험을 하나의 일관된 해석
  구조로 묶어 설명할 수 있다.

  \item[Level C: 경험 과학]
  이 프레임워크가 과학 이론이 되려면 관측 가능한 예측과 반증 가능성이
  필요하다. 본 논문은 그 단계에 도달했다고 주장하지 않는다.
\end{description}

따라서 ``재민 우주론''은 검증된 사실의 체계가 아니라, 고차원 정보 현실과
인지적 투영 사이의 관계를 탐구하기 위한 형식적 연구 프로그램으로 이해되어야
한다.

\section{동기 부여 비유: 이차원 개미 (Motivating Analogy: The Two-Dimensional Ant)}

이차원 평면 위에서만 움직이고 감각하는 개미를 생각하자. 개미는 평면
내부의 길이, 방향, 경계를 지각할 수 있지만, 평면 밖의 세 번째 축을 직접
경험하지 못한다. 삼차원 관찰자는 개미가 이해하지 못하는 방식으로 평면을
구부리거나, 평면 밖에서 개미의 위치를 관찰하거나, 개미에게는 갑작스럽게
보이는 변화를 만들 수 있다.

이 비유는 인간이 실제로 이차원 개미와 같다는 주장이 아니다. 핵심은
지각 가능한 차원이 존재론적 차원 전체와 동일하다는 가정이 필수적이지
않다는 점이다. 인간의 삼차원 의식 경험도 더 높은 차원 정보 상태의
인지적 투영일 수 있다. 이러한 가능성을 형식화하려면 전체 상태공간,
의식 투영, 잠재의식 표현, 시간 매개변수를 구별해야 한다.

\section{핵심 공리 (Core Postulates)}

\begin{axiom}[고차원 현실]
현실의 실제 상태공간은 $n>3$인 $n$차원 정보 다양체 $M^n$이다.
\end{axiom}

\begin{axiom}[의식 투영]
인간의 의식 경험은 전체 상태 $X$를 삼차원 인지 인터페이스로 사상하는
저차원 투영
\begin{equation}
  C = P(X), \qquad P:M^n \to \mathbb{R}^3
\end{equation}
으로 표현된다.
\end{axiom}

\begin{axiom}[잠재의식의 고차원 처리]
잠재의식 표현은
\begin{equation}
  S = Q(X), \qquad Q:M^n \to \mathbb{R}^m,\qquad 3<m\le n
\end{equation}
으로 주어지며, 의식보다 더 높은 차원의 정보 표현을 포함할 수 있다.
\end{axiom}

\begin{axiom}[뇌-상위차원 인터페이스]
뇌는 고차원 정보공간과 삼차원 의식 경험 사이의 인터페이스로 모델링될 수
있다. 신경 상태를 $N_b$, 고차원 정보장을 $H$, 접근 또는 결합 효율을
$\eta\in[0,1]$로 두면 잠재의식 표현은 더 일반적으로
\begin{equation}
  S = Q_\eta(X,N_b,H)
\end{equation}
로 표현된다.
\end{axiom}

\begin{axiom}[직관의 압축]
직관은 잠재의식의 고차원 표현에서 의식 수준으로 압축된 신호
\begin{equation}
  I = G(S), \qquad G:\mathbb{R}^m \to \mathbb{R}^k
\end{equation}
이다.
\end{axiom}

\begin{axiom}[가변 내부 시간]
내부 우주 시간 $\tau$와 고차원 기준 시간 $T$는
\begin{equation}
  \frac{d\tau}{dT} = \alpha(X,T), \qquad \alpha(X,T)>0
\end{equation}
으로 연결된다.
\end{axiom}

\begin{axiom}[내부 불변성]
내부 관찰자는 모든 내부 측정 장치와 물리 과정이 $\tau$에 의해 기술되므로,
모든 내부 과정이 같은 비율로 스케일되는 경우 그 스케일 인자 $\alpha$를
내부 시계만으로 직접 검출할 수 없다.
\end{axiom}

\begin{axiom}[약한 잠재의식 민감도]
잠재의식은 $\alpha$ 또는 $d\alpha/dT$의 변화를 약하게 부호화할 수 있으며,
그 결과 주관적 시간 가속감이 형성될 수 있다.
\end{axiom}

\begin{axiom}[지능의 인터페이스 효율]
지능, 통찰, 창의성은 신경 계산 능력만의 함수가 아니라 신경 연결 복잡도와
고차원 정보 접근 효율 $\eta$의 결합 효과로 해석될 수 있다. 이 공리는
검증된 신경과학 명제가 아니라 본 프레임워크의 해석적 가정이다.
\end{axiom}

\section{수학적 틀 (Mathematical Framework)}

\subsection{기본 상태공간과 표기}

고차원 현실은 정보 다양체 $M^n$으로 표현한다. 고차원 기준 시간 $T$에서
우주의 실제 상태는
\begin{equation}
  X(T) \in M^n,\qquad n>3
\end{equation}
이다. 인간 의식이 접근하는 상태는 $X(T)$ 자체가 아니라 투영된 인터페이스
$C=P(X)$이다.

\begin{center}
\begin{tabular}{>{\raggedright\arraybackslash}p{0.23\linewidth}p{0.67\linewidth}}
\toprule
\textbf{기호} & \textbf{의미} \\
\midrule
$M^n$ & $n>3$인 고차원 정보 다양체 \\
$X(T)$ & 고차원 기준 시간 $T$에서의 실제 우주 상태 \\
$T$ & 고차원 기준 시간 \\
$\tau$ & 내부 관찰자가 사용하는 내부 우주 시간 \\
$H$ & 고차원 정보장 또는 정보 구조 \\
$\mathcal{H}$ & 가능한 고차원 정보장들의 공간 \\
$N_b$ & 뇌의 내부 신경 상태 \\
$\mathcal{N}$ & 가능한 신경 상태들의 공간 \\
$\Gamma$ & 신경 연결 구조 또는 연결 복잡도 \\
$\eta$ & 고차원 정보 접근 또는 결합 효율, $0\le\eta\le 1$ \\
$P$ & $M^n$에서 $\mathbb{R}^3$로 가는 의식 투영 사상 \\
$C$ & 의식적 삼차원 인터페이스, $C=P(X)$ \\
$Q$ & $M^n$에서 $\mathbb{R}^m$으로 가는 잠재의식 표현 사상 \\
$Q_\eta$ & 신경 상태와 고차원 정보장의 결합을 포함하는 잠재표현 사상 \\
$S$ & 잠재의식 고차원 표현, 기본형 $S=Q(X)$ 또는 확장형 $S=Q_\eta(X,N_b,H)$ \\
$G$ & 잠재의식 표현을 직관 신호로 압축하는 사상 \\
$I$ & 직관 신호, $I=G(S)$ \\
$\alpha(X,T)$ & 내부 시간과 기준 시간을 연결하는 양의 시간율 함수 \\
$Y$ & 예측하거나 판단하려는 사건 또는 의사결정 대상 \\
\bottomrule
\end{tabular}
\end{center}

\begin{definition}[재민 우주론의 최소 구조]
재민 우주론의 최소 수학 구조는
\begin{equation}
  \mathcal{J} = \left(M^n, X, H, N_b, \eta, P, Q_\eta, G, \alpha\right)
\end{equation}
의 요소로 구성된다. 여기서 $M^n$은 고차원 상태공간, $X$는 실제 상태 궤적,
$H$는 고차원 정보장, $N_b$는 신경 상태, $\eta$는 접근 효율, $P$는 의식
투영, $Q_\eta$는 결합 효율을 포함한 잠재의식 표현, $G$는 직관 압축,
$\alpha$는 내부 시간율 함수이다.
\end{definition}

\subsection{간단한 장난감 모형}

가장 단순한 예로
\begin{equation}
  M^4=\mathbb{R}^4,\qquad X=(x,y,z,w)
\end{equation}
를 둔다. 의식 투영은 네 번째 좌표 $w$를 버리는 사상
\begin{equation}
  P(x,y,z,w)=(x,y,z)
\end{equation}
이고, 잠재의식 표현은
\begin{equation}
  Q_\eta(x,y,z,w,N_b,H)=(x,y,z,\eta\epsilon w+\xi)
\end{equation}
라고 하자. 여기서 $\epsilon$은 네 번째 좌표의 신호 강도, $\eta$는
고차원 정보 접근 효율, $\xi$는 잡음이다.

사건 $Y$가 $w$에 부분적으로 의존한다면 의식 표현 $C=(x,y,z)$만으로는
$Y$를 완전히 예측할 수 없다. 그러나 $\eta\epsilon>0$이고 잡음 $\xi$가 충분히
제한되어 있다면, 잠재의식 표현 $S=Q_\eta(X,N_b,H)$는 $Y$에 관한 추가 정보를 보존할
수 있다. 이 경우 직관 $I=G(S)$는 완전한 예측이 아니라 불확실성을 줄이는
압축 신호로 작동한다.

\section{저차원 투영으로서의 의식 (Consciousness as Low-Dimensional Projection)}

의식은 현실 전체의 사본이 아니라, 인지적으로 조작 가능한 인터페이스로
정의된다. 핵심 방정식은
\begin{equation}
  C=P(X), \qquad P:M^n\to\mathbb{R}^3
\end{equation}
이다. 여기서 $P$는 전체 정보 상태에서 행동, 지각, 공간적 탐색에 필요한
일부 구조만을 추출하는 사상이다.

\begin{remark}
$P$는 반드시 선형 투영일 필요가 없다. 비선형 사상, 확률적 사상,
정보 병목을 포함하는 사상도 허용된다. 중요한 점은 $P$가 차원을 낮추며,
그 과정에서 일부 정보가 의식 수준에서 사라질 수 있다는 것이다.
\end{remark}

의식이 삼차원 공간을 안정적으로 경험한다는 사실은 $C$가
$\mathbb{R}^3$의 좌표 구조를 가진다는 모델링 선택으로 표현된다. 이 선택은
인간 경험의 형식을 수학적으로 반영할 뿐, 실제 세계가 오직 삼차원이라는
결론을 전제하지 않는다.

\section{상위차원 인터페이스로서의 뇌 (The Brain as a Higher-Dimensional Interface)}

재민 우주론의 중심축은 시간 가속 자체가 아니라
뇌의 역할이다. 이 관점에서 뇌는 모든 계산을 삼차원 내부에서 닫힌 방식으로
수행하는 장치가 아니라, 고차원 정보장 $H$와 의식 인터페이스 $C$ 사이의
변환기이다. 이를 일반식으로 쓰면
\begin{equation}
  S = Q_\eta(X,N_b,H)
  \label{eq:q-eta}
\end{equation}
이다. 여기서 $N_b$는 뇌의 물리적 신경 상태, $H$는 고차원 정보장,
$\eta$는 고차원 정보 접근 효율이다.

신경과학에서 말하는 고차원 신경 상태공간은 통상 수학적 표현 공간으로
해석된다. 재민 우주론은 이 해석을 부정하지 않는다. 다만 추가 공리로,
그러한 고차원 상태공간이 실제 고차원 정보 구조 $M^n$ 또는 $H$와 부분적으로
대응할 수 있다고 가정한다.

\begin{assumption}[고차원 신경 표현의 존재론적 해석]
인지계에서 관찰되거나 모델링되는 고차원 신경 상태공간은 단순한 수학적
기술일 수 있다. 그러나 재민 우주론에서는 그 상태공간이 고차원 정보 다양체
$M^n$ 안의 구조와 부분적으로 결합될 수 있다고 가정한다.
\end{assumption}

이 가정은 뇌 가소성이나 연결 복잡도에 관한 기존 설명을 대체하지 않는다.
오히려 다음과 같은 이중 구조를 둔다.
\begin{equation}
  S = (1-\eta)Q_N(N_b) + \eta Q_H(X,H) + \xi,
  \label{eq:interface-mixture}
\end{equation}
여기서 $Q_N$은 통상적 신경 계산 성분, $Q_H$는 고차원 정보 결합 성분,
$\xi$는 잡음 또는 왜곡이다. $\eta=0$이면 표준적 신경 계산 모델에 가까워지고,
$\eta>0$이면 고차원 정보 결합 항이 잠재의식 표현에 기여한다.

\section{잠재의식 고차원 표현 (Subconscious High-Dimensional Representation)}

잠재의식은 다음과 같이 더 높은 차원의 표현을 가질 수 있다.
\begin{equation}
  S=Q_\eta(X,N_b,H), \qquad Q_\eta:M^n\times\mathcal{N}\times\mathcal{H}\to\mathbb{R}^m,
  \qquad 3<m\le n.
\end{equation}
이 표현은 초자연적 능력을 뜻하지 않는다. 현대 인지과학과 신경과학에서도
신경 상태공간, 잠재 변수, 표현 학습, 예측 처리와 같은 개념은 고차원
상태표현을 사용한다. 재민 우주론은 이러한 고차원 신경 상태공간을 더
넓은 고차원 정보 구조에 대한 인지적 인터페이스로 해석할 가능성을 둔다.

\begin{assumption}[정보 보존성]
어떤 사건 또는 판단 대상 $Y$에 대해, $Q$는 $P$가 잃는 관련 정보를 일부
보존할 수 있다.
\end{assumption}

이 가정은 모든 상황에서 참이라고 주장되지 않는다. 잠재의식 표현도 잡음,
편향, 학습 오류, 기억 왜곡을 포함할 수 있다.

\section{압축된 고차원 정보로서의 직관 (Intuition as Compressed High-Dimensional Information)}

직관은
\begin{equation}
  I=G(S), \qquad G:\mathbb{R}^m\to\mathbb{R}^k
\end{equation}
로 정의된다. 여기서 $k$는 의식 수준에서 처리 가능한 신호 차원이다.
직관은 완전한 예언이나 오류 없는 판단이 아니다. 오히려 고차원 잠재표현이
의식으로 전달될 때 나타나는 압축되고 손실 있는 요약이다.

정보이론적으로, 특정 사건 또는 의사결정 대상 $Y$에 대해 잠재의식 표현이
의식 표현보다 더 많은 관련 정보를 보존한다면
\begin{equation}
  I(Y;S)\ge I(Y;C)
\end{equation}
이고
\begin{equation}
  H(Y\mid S)\le H(Y\mid C)
\end{equation}
일 수 있다. 그러나 불확실성은 사라지지 않는다.
\begin{equation}
  H(Y\mid S)>0.
\end{equation}

\begin{proposition}[잠재표현의 불확실성 감소]
$\dim(S)>\dim(C)$이고, $S$가 투영 $P$에서 손실되는 $Y$ 관련 정보를 보존한다면,
잠재의식 표현은 의식 표현보다 조건부 불확실성을 줄일 수 있다.
\begin{equation}
  H(Y\mid S)\le H(Y\mid C).
\end{equation}
\end{proposition}

\begin{proof}
$S$가 $C$보다 $Y$와 관련된 정보를 더 많이 보존한다는 가정은
$I(Y;S)\ge I(Y;C)$로 표현된다. 엔트로피 항등식
$I(Y;Z)=H(Y)-H(Y\mid Z)$를 $Z=S$와 $Z=C$에 각각 적용하면,
$H(Y)-H(Y\mid S)\ge H(Y)-H(Y\mid C)$가 된다. 따라서
$H(Y\mid S)\le H(Y\mid C)$이다. 다만 이 부등식은
$H(Y\mid S)=0$을 의미하지 않으므로, 직관은 오류 가능성을 계속 가진다.
\end{proof}

\section{지능과 인터페이스 효율 (Intelligence and Interface Efficiency)}

재민 우주론에서 지능 차이는 단순히 뉴런 수의 차이로 환원되지 않는다.
기존 신경과학적 설명은 연결 구조, 연결 복잡도, 정보 흐름, 학습 경험,
뇌 가소성 같은 요인을 강조한다. 본 프레임워크는 이를 유지하면서, 추가로
고차원 정보 접근 효율 $\eta$를 둔다.

신경 연결 구조를 $\Gamma$라고 하자. 인지적 성과 또는 추상적 지능 점수를
$\mathcal{Z}$로 놓으면, 가장 단순한 형식은
\begin{equation}
  \mathcal{Z}
  =
  \phi(N_b,\Gamma) + \eta\,\psi(S)
  \label{eq:intelligence}
\end{equation}
이다. 여기서 $\phi$는 통상적 신경 계산 능력, $\psi$는 고차원 잠재표현이
문제 해결에 기여하는 항이다. 이 식은 실제 IQ 측정식을 주장하지 않는다.
오직 ``신경 계산 능력''과 ``고차원 정보 결합 효율''을 분리하기 위한
형식적 모형이다.

\begin{proposition}[동일한 신경 자원에서도 다른 인지 성과가 가능함]
두 인지계가 동일한 $N_b$와 $\Gamma$를 가진다고 가정하더라도,
식~\eqref{eq:intelligence}에서 $\eta$가 다르면 $\mathcal{Z}$가 달라질 수
있다.
\end{proposition}

\begin{proof}
$N_b$와 $\Gamma$가 같으면 $\phi(N_b,\Gamma)$ 항은 동일하다. 그러나
$\psi(S)$가 0이 아니고 $\eta$ 값이 서로 다르면 $\eta\psi(S)$ 항이 달라진다.
따라서 전체 $\mathcal{Z}$도 달라질 수 있다.
\end{proof}

이 명제는 경험적 사실을 확정하지 않는다. 다만 비슷한 물리적 신경 자원을
가진 개체 사이에서도 통찰, 직관, 창의성의 차이가 발생할 수 있다는 가설을
수학적으로 표현한다.

\section{AI와 저차원 계산 부담 (AI and Low-Dimensional Computational Burden)}

또 다른 해석적 응용은 인간과 현재 AI의 계산 방식 차이다.
현재 AI 시스템은 고차원 벡터 표현을 사용하지만, 그 연산은 물리적으로는
저차원 우주 내부의 하드웨어에서 수행된다. 재민 우주론의 언어로 표현하면
대부분의 현재 AI는
\begin{equation}
  \eta_{\mathrm{AI}} = 0
\end{equation}
또는 고차원 정보장 $H$와의 직접 결합을 가정하지 않는 체계로 모델링된다.
따라서 고차원 문제를 해결하려면 매개변수, 메모리, 전력, 병렬 하드웨어를
통해 내부적으로 근사해야 한다.

반대로 인간 뇌에 대해 $\eta_{\mathrm{human}}>0$이라는 가설을 둔다면,
인간의 낮은 에너지 사용량과 강한 직관적 문제 해결 능력을 ``계산량''만이
아니라 ``접근 효율''의 문제로 해석할 수 있다. 이 해석은 현재 AI 연구나
신경과학의 확정된 결론이 아니다. 논문 수준에서는 다음과 같은 조건부 명제로
남겨야 한다.

\begin{quote}
만약 인간 뇌가 고차원 정보장과 결합하는 인터페이스라면, 인간 지능의 일부는
삼차원 내부 계산량이 아니라 고차원 정보 접근 효율 $\eta$에 의해 설명될 수
있다.
\end{quote}

\section{가변 내부 시간 (Variable Internal Time)}

고차원 기준 시간 $T$와 내부 우주 시간 $\tau$를 구별한다. 두 시간은
\begin{equation}
  \frac{d\tau}{dT}=\alpha(X,T),\qquad \alpha(X,T)>0
  \label{eq:time-rate}
\end{equation}
으로 연결된다. 만약 어떤 구간에서
\begin{equation}
  \frac{d\alpha}{dT}>0
\end{equation}
라면 고차원 관찰자는 내부 우주의 시간이 점점 빠르게 진행된다고 묘사할 수
있다. 그러나 내부 관찰자는 자신의 물리 법칙을 $T$가 아니라 $\tau$에 대해
측정한다.

내부 물리 상태 $x$가
\begin{equation}
  \frac{dx}{d\tau}=f(x)
\end{equation}
를 따른다면, 고차원 기준 시간에 대해서는 연쇄법칙에 의해
\begin{equation}
  \frac{dx}{dT}
  =
  \frac{d\tau}{dT}\frac{dx}{d\tau}
  =
  \alpha(X,T)f(x)
  \label{eq:external-dynamics}
\end{equation}
이다.

\begin{proposition}[내부 시계만으로는 균일한 시간율을 직접 검출할 수 없음]
모든 내부 물리 과정이 $\tau$에 대해 진화하고
$d\tau/dT=\alpha(T)$라면, 내부 관찰자는 내부 시계만으로 $\alpha$를 직접
검출할 수 없다.
\end{proposition}

\begin{proof}
내부의 시계, 신경 과정, 화학 과정, 측정 표준은 모두 $\tau$를 기준으로
진화한다. 고차원 기준 시간에서 보면 각 과정은 식~\eqref{eq:external-dynamics}처럼
같은 시간율 인자 $\alpha$를 받는다. 내부 관찰자가 측정하는 것은 과정들
사이의 비율과 상관관계인데, 모든 관련 과정이 동일한 내부 시간 $\tau$로
매개화되면 이러한 내부 비율은 보존된다. 따라서 균일한 $\alpha$ 스케일링은
내부 시계만으로 직접 검출되지 않는다.
\end{proof}

\section{노화와 주관적 시간 가속 (Subjective Time Acceleration with Aging)}

나이가 들수록 시간이 빠르게 흐르는 것처럼 느끼는 현상은 일반적으로 기억
밀도, 새로움의 감소, 반복 경험의 증가, 주의 자원의 변화 등으로 설명된다.
재민 우주론은 이러한 심리학적 설명을 배제하지 않고, 약한 고차원 시간율
민감도 항을 추가하는 방식으로 모델링한다.

내부 시간 간격 $\Delta\tau$에 대한 체감 지속시간을
\begin{equation}
  D_{\mathrm{felt}}(a,T)
  =
  \mu(a)\,[\alpha(T)]^{-\lambda}\,\Delta\tau
  \label{eq:felt-duration}
\end{equation}
로 둔다. 여기서 $a$는 생물학적 나이, $\mu(a)$는 기억 밀도 또는 새로움
밀도 함수, $\lambda\ge 0$는 고차원 내부 시간율 변화에 대한 잠재의식의
민감도이다.

\begin{assumption}[기억 밀도의 감소]
생물학적 나이가 증가할수록 평균적 새로움 또는 기억 밀도는 감소한다.
\begin{equation}
  \frac{d\mu}{da}<0.
\end{equation}
\end{assumption}

특수한 경우 $\lambda=0$이면 식~\eqref{eq:felt-duration}은 표준적 심리
모형만을 나타낸다. 반대로 $\lambda>0$이고 $\alpha(T)$가 증가한다면,
재민 우주론은 약한 고차원 시간율 민감도가 체감 지속시간을 추가로 감소시킬
수 있다고 해석한다.

\begin{proposition}[체감 지속시간의 감소 조건]
$d\mu/da<0$, $\lambda\ge 0$이고 $\alpha(T)$가 비감소 함수라면,
식~\eqref{eq:felt-duration}의 체감 지속시간은 매개변수 값에 따라 나이와
우주 기준 시간에 대해 감소하는 경향을 가진다.
\end{proposition}

\begin{proof}
$D_{\mathrm{felt}}$는 $\mu(a)$에 비례하므로 $d\mu/da<0$이면 나이 증가에
따라 감소하는 성분을 가진다. 또한 $\lambda\ge 0$이고 $\alpha(T)$가
비감소이면 $[\alpha(T)]^{-\lambda}$는 비증가 함수이다. 두 항이 모두
양수이므로, 다른 조건이 일정할 때 두 효과는 체감 지속시간을 감소시키는
방향으로 작용한다. 실제 감소율은 $\mu$, $\alpha$, $\lambda$의 구체적
형태에 의존한다.
\end{proof}

\section{수명 기대와 문명 가속 (Life Expectancy and Civilizational Acceleration)}

이 프레임워크는 내부 시간의 가속이 인간 수명을 직접 증가시킨다고 주장하지
않는다. 수명 기대는 내부 시간 $\tau$에서 측정되는 생물학적, 의학적,
사회적 변수의 결과이다. 가능한 간접 경로는 문명과 생의학 기술의 발전이다.

생의학 기술 수준을 $B(\tau)$라고 하자. 내부 시간 기준 수명 기대를
\begin{equation}
  L_\tau=L_0+\beta B(\tau),\qquad \beta>0
  \label{eq:life}
\end{equation}
로 모델링한다. 여기서 $L_0$는 기초 생물학적 수명이고, $\beta$는 생의학
기술이 수명 기대에 기여하는 정도이다.

기술 발전은 고차원 잠재정보 처리 효율 $E(S,\eta)$에 부분적으로 의존할 수 있다.
\begin{equation}
  \frac{dB}{d\tau}=\rho E(S,\eta)B,\qquad \rho>0.
  \label{eq:biomed}
\end{equation}
예시적으로
\begin{equation}
  E(S,\eta)=E_0+\gamma\eta\log(1+\alpha(T)),\qquad \gamma\ge 0
\end{equation}
또는 더 일반적으로
\begin{equation}
  E(S,T,\eta)=E_0+\gamma\Phi(S,\alpha,\eta)
\end{equation}
를 둘 수 있다. 이 해석에서 고차원 정보 처리는 창의성, 직관, 과학적 발견,
의학적 진보의 속도에 간접적으로 기여할 수 있다. 하지만 이는 가능한 수학적
경로일 뿐, 현재 경험적으로 확정된 인과관계가 아니다.

\section{필수 회전 없는 일반화 모형 (Generalized Model Without Mandatory Rotation)}

재민 우주론의 기본 시간율 방정식은
\begin{equation}
  \frac{d\tau}{dT}=\alpha(X,T)
\end{equation}
이다. 이 식은 우주가 고차원에서 반드시 회전해야 한다는 공리를 포함하지
않는다. $\alpha$를 증가시키는 가능한 메커니즘은 다양할 수 있다.

\begin{itemize}[leftmargin=2em]
  \item 고차원 스케일 팽창
  \item 곡률 변화
  \item 정보 밀도 변화
  \item 고차원 다양체 안에서의 브레인 유사 운동
  \item 고차원 매장공간에서의 회전
\end{itemize}

따라서 회전은 하나의 가능한 생성 메커니즘이지, 프레임워크 전체의 필수
전제가 아니다.

\section{선택적 회전 메커니즘 (Optional Rotational Mechanism as a Special Case)}

회전 메커니즘을 선택하는 경우, 고차원 중심 $O$와 반지름
\begin{equation}
  r(T)=\|X(T)-O\|
\end{equation}
를 정의할 수 있다. 각속도를
\begin{equation}
  \frac{d\theta}{dT}=\Omega
\end{equation}
로 두면 접선 속도는
\begin{equation}
  v(T)=\Omega r(T)
\end{equation}
이다. 시간율 함수는
\begin{equation}
  \alpha(T)=F(v(T)),\qquad F'(v)>0
\end{equation}
처럼 둘 수 있다. 예를 들어
\begin{equation}
  \alpha(T)=1+\kappa v(T)^2,\qquad \kappa>0
\end{equation}
는 속도가 증가할수록 내부 시간율이 증가하는 간단한 모형이다.

\begin{proposition}[회전은 충분조건이지만 필요조건은 아님]
위와 같은 회전 모형은 증가하는 $\alpha(T)$를 생성할 수 있으므로 충분조건의
예가 된다. 그러나 재민 우주론의 기본 방정식은 회전을 요구하지 않으므로,
회전은 필요조건이 아니다.
\end{proposition}

\begin{proof}
$F'(v)>0$이면 $v(T)$가 증가하는 구간에서 $\alpha(T)=F(v(T))$도 증가한다.
따라서 회전 반지름 또는 접선 속도가 증가하는 회전 모형은 시간율 증가를
생성할 수 있다. 반면 기본 방정식 $d\tau/dT=\alpha(X,T)$는 $\alpha$의
원인을 특정하지 않으므로, 곡률 변화나 정보 밀도 변화 같은 다른 원인도
형식적으로 허용된다. 그러므로 회전은 충분하지만 필요하지 않다.
\end{proof}

\section{논의 (Discussion)}

본 프레임워크의 장점은 의식, 잠재의식, 직관, 지능 차이, 시간 경험을 하나의
수학 구조 안에서 연결한다는 점이다. 의식은 $P$를 통한 저차원 투영이고,
잠재의식은 $Q_\eta$를 통한 고차원 표현이며, 직관은 $G$에 의한 압축 신호이다.
뇌는 $N_b$와 $H$를 연결하는 인터페이스로 모델링되고, 접근 효율 $\eta$는
통찰과 인지 성과의 차이를 설명하기 위한 가설적 매개변수로 사용된다.
내부 시간은 $\tau$, 고차원 기준 시간은 $T$로 분리되며, 두 시간의 관계는
$\alpha(X,T)$에 의해 조절된다.

이 구조에서 시간 가속 가설은 중심 공리가 아니라 파생적 응용에 가깝다.
처음에는 내부 시간율 변화가 문제를 제기하지만, 더 일반적인 형식화는
``뇌가 고차원 정보공간과 결합하는가''라는 질문에 놓인다. 만약 이 결합이
가능하다면 직관, 창의성, 지능 차이, 주관적 시간 감각은 같은 인터페이스
모형의 서로 다른 표현으로 해석될 수 있다.

그러나 이러한 통합성은 경험적 타당성을 보장하지 않는다. 특히 균일한
시간율 스케일링은 내부 관찰자에게 재매개화와 구별되지 않을 수 있다.
따라서 이 프레임워크가 과학 이론으로 발전하려면 $\alpha$, $S$, $I$,
$\eta$에 관한 독립적 측정 가능량이나, 표준 심리학 및 신경과학 모형과
구별되는 예측이 필요하다.

\section{한계 (Limitations)}

\begin{enumerate}[leftmargin=2em]
  \item 현재 이 프레임워크에 대해 제공된 경험적 증거는 없다.
  \item $\alpha(X,T)$를 강제하는 특정 물리 메커니즘이 지정되어 있지 않다.
  \item 신경 고차원 상태공간과 실제 고차원 다양체 $M^n$ 사이의 대응이
  불분명하다.
  \item 고차원 정보 접근 효율 $\eta$를 직접 측정하는 방법이 아직 없다.
  \item 고차원 기준 시간 $T$를 직접 측정하는 방법이 제시되어 있지 않다.
  \item 균일한 내부 시간 가속은 단순한 시간 재매개화와 관측상 동등할 수
  있다.
  \item 검증 가능한 예측이 없으면 이 틀은 반증 불가능성의 위험을 가진다.
  \item 직관은 고차원 우주론 없이도 일반적인 잠재의식 패턴 인식만으로
  설명될 수 있다.
  \item 인간 뇌가 현재 AI보다 효율적인 이유는 기존 신경과학과 컴퓨터
  공학만으로도 부분적으로 설명될 수 있으므로, 고차원 인터페이스 가설은
  별도의 경험적 근거를 요구한다.
  \item 극단적 뇌 손상 후 기능 보존 사례는 신중한 임상 해석이 필요하며,
  이를 곧바로 고차원 정보 결합의 증거로 사용할 수 없다.
\end{enumerate}

\section{향후 연구 (Future Work)}

\begin{enumerate}[leftmargin=2em]
  \item $P$, $Q$, $G$의 구체적 함수 형태를 형식화한다.
  \item $Q_\eta$와 접근 효율 $\eta$의 측정 가능한 대리 변수를 정의한다.
  \item 잠재의식 표현 $S$의 측정 가능한 대리 변수를 정의한다.
  \item 이 모형을 신경 잠재공간 및 표현 학습 모델과 비교한다.
  \item 연결 복잡도 $\Gamma$, 뇌 가소성, 고차원 표현 사이의 관계를
  수학적으로 모델링한다.
  \item 인간과 AI의 에너지 효율 차이를 $\eta=0$ 및 $\eta>0$ 조건부 모형으로
  비교한다.
  \item 정보 병목 이론과 직관 압축 $G$의 관계를 연구한다.
  \item $M^n$의 미분기하학적 모형을 탐색한다.
  \item 주관적 시간 지각 자료가 $\lambda$를 제약할 수 있는지 조사한다.
  \item 후보 경험 예측을 도출한다.
  \item 시뮬레이션 모델을 개발한다.
  \item $n=4$, $m=4$, $C\subset\mathbb{R}^3$인 장난감 모형을 확장한다.
\end{enumerate}

\section{결론 (Conclusion)}

재민 우주론은 인간 의식이 고차원 정보 현실의 저차원 투영일 수 있다는
가정을 수학적으로 정리한 사변적 프레임워크이다. 이 틀은 $M^n$, $X(T)$,
$H$, $N_b$, $\eta$, $P$, $Q_\eta$, $G$, $\alpha$를 중심으로 의식,
잠재의식, 직관, 지능, 내부 시간, 주관적 시간 가속, 문명적 발전을 연결한다.
후속 형식화에서 정리되는 가장 중요한 수정점은 시간 가속이 이론의 출발점일 수는
있지만, 이론의 중심은 뇌가 상위차원 정보공간의 인터페이스라는 명제에
있다는 것이다. 핵심 결과 중 하나는 내부 관찰자가 내부 시간 $\tau$에 의해
모든 물리 과정을 측정한다면, 고차원 기준 시간 $T$에서의 균일한 시간율
변화가 내부적으로 직접 검출되지 않을 수 있다는 점이다.

그럼에도 이 프레임워크는 현재 검증된 물리 이론이 아니며, 경험적 확증을
주장하지 않는다. 그것은 수학적으로 정의 가능한 해석적 모형이다. 향후
연구의 핵심 과제는 이 모형이 표준 심리학, 신경과학, 물리학적 설명과
구별되는 측정 가능한 예측을 생성할 수 있는지 밝히는 것이다.

\nocite{*}
\bibliographystyle{plain}
\bibliography{references}

\end{document}
